\chapter{The Reciprocal Paired-Spectrum Polynomial}
\label{chap:reciprocal-paired-spectrum-polynomial}

\status{Exact finite-set algebraic reduction; no claim that the paired spectrum is nonempty and no resolution of $F(-1/4)$.}

\section{Purpose}

Chapters~\ref{chap:quotient-negative-inversion-zero-set-no-go}--\ref{chap:v4-paired-spectrum-orbits} show that the quotient paired-root set
\[
 P_J:=\left\{w:F(w)=0,\ F\!\left(\frac1{16w}\right)=0\right\}
\]
is finite and invariant under
\[
 J(w)=\frac1{16w}.
\]
The present chapter packages this finite set into a canonical monic polynomial and derives its exact reciprocal symmetry.

\section{Normalized reciprocal coordinate}

Define
\[
 \lambda=4w.
\]
Then
\[
 4J(w)=\frac4{16w}=\frac1{4w}=\lambda^{-1}.
\]
Thus the quotient negative inversion is conjugated to the standard reciprocal map
\[
 \boxed{\lambda\mapsto\lambda^{-1}}.
\]

Let
\[
 \Lambda:=4P_J.
\]
Since $F(0)=\xi(1/2)>0$, the paired set contains no $w=0$, so
\[
 \Lambda\subset\mathbb C^\times.
\]

\section{Reciprocal and conjugation closure}

\begin{proposition}
The finite set $\Lambda$ is invariant under reciprocal inversion:
\[
 \lambda\in\Lambda\Longrightarrow\lambda^{-1}\in\Lambda.
\]
\end{proposition}

\begin{proof}
If $\lambda=4w$ with $w\in P_J$, then $J(w)\in P_J$.  Therefore
\[
 4J(w)=\lambda^{-1}\in\Lambda.
\]
\end{proof}

The quotient entire function $F$ has real Taylor coefficients because
\[
 \xi\!\left(\frac12+z\right)=F(z^2)
\]
and the centered xi function is real on the real axis. Hence
\[
 F(\bar w)=\overline{F(w)}.
\]
Since $J$ has real coefficients,
\[
 J(\bar w)=\overline{J(w)}.
\]
Thus $P_J$, and therefore $\Lambda$, is also invariant under complex conjugation.

\section{Canonical monic polynomial}

Let
\[
 n:=|P_J|=|\Lambda|
\]
and define
\[
 Q(x):=\prod_{\lambda\in\Lambda}(x-\lambda).
\]
If the paired set is empty, the empty-product convention gives $Q(x)=1$.

By Chapter~\ref{chap:v4-paired-spectrum-orbits},
\[
 n=2a+\varepsilon,
 \qquad
 \varepsilon\in\{0,1\},
\]
where $\varepsilon=1$ exactly when the exceptional reciprocal fixed point
\[
 \lambda=-1
\]
is present.

The other reciprocal fixed point
\[
 \lambda=+1
\]
corresponds to $w=+1/4$ and is excluded by the positive-coefficient theorem,
\[
 F(1/4)>0.
\]

\section{Unit constant term}

Every nonfixed reciprocal orbit has the form
\[
 \{\lambda,\lambda^{-1}\}
\]
and therefore has product $1$. The only possible fixed contribution is $-1$, so
\[
 \prod_{\lambda\in\Lambda}\lambda=(-1)^\varepsilon.
\]
Because $n\equiv\varepsilon\pmod2$,
\[
 Q(0)
 =(-1)^n\prod_{\lambda\in\Lambda}\lambda
 =1.
\]
Hence
\[
 \boxed{Q(0)=1}.
\]

\section{Self-reciprocal functional identity}

\begin{theorem}[SOH-G018 reciprocal paired-spectrum polynomial]
The monic polynomial $Q$ satisfies
\[
 \boxed{Q(x)=x^nQ(1/x)}.
\]
\end{theorem}

\begin{proof}
We have
\[
 x^nQ(1/x)
 =x^n\prod_{\lambda\in\Lambda}\left(\frac1x-\lambda\right)
 =\prod_{\lambda\in\Lambda}(1-\lambda x).
\]
Its roots are precisely the reciprocal points $\lambda^{-1}$. Reciprocal closure of $\Lambda$ gives the same root set as $Q$. Its leading coefficient is
\[
 (-1)^n\prod_{\lambda\in\Lambda}\lambda
 =Q(0)
 =1.
\]
Therefore both sides are monic degree-$n$ polynomials with the same roots, hence they are identical.
\end{proof}

\section{Palindromic real coefficients}

Write
\[
 Q(x)=x^n+c_1x^{n-1}+\cdots+c_{n-1}x+1.
\]
The self-reciprocal identity implies
\[
 \boxed{c_k=c_{n-k}}.
\]
Thus the coefficient vector is palindromic.

Conjugation closure of $\Lambda$ implies that the multiset of roots of $Q$ is unchanged by conjugation. Since $Q$ is monic,
\[
 \overline{Q(\bar x)}=Q(x),
\]
so
\[
 \boxed{Q\in\mathbb R[x]}.
\]

\section{Orbit factorization}

Choose one representative $\lambda_j$ from each nonfixed reciprocal two-cycle and define
\[
 \tau_j:=\lambda_j+\lambda_j^{-1}.
\]
Then
\[
 (x-\lambda_j)(x-\lambda_j^{-1})
 =x^2-\tau_jx+1.
\]
Therefore
\[
 \boxed{
 Q(x)
 =(x+1)^\varepsilon
 \prod_{j=1}^{a}\left(x^2-\tau_jx+1\right).
 }
\]
Individual $\tau_j$ need not be real, but the collection is closed under conjugation and the complete product has real coefficients.

\section{Fixed-point criterion}

Since $\lambda=+1$ is excluded,
\[
 \boxed{Q(1)\neq0}.
\]
The only possible reciprocal fixed root is $\lambda=-1$, corresponding to $w=-1/4$. Hence
\[
 \boxed{
 Q(-1)=0
 \iff
 \varepsilon=1
 \iff
 F(-1/4)=0.
 }
\]
This is an exact equivalence, not a resolution of the final condition.

\section{Equivalent identity in the original quotient coordinate}

Define
\[
 H(w):=\prod_{r\in P_J}(w-r).
\]
Because
\[
 Q(x)=4^nH(x/4),
\]
the self-reciprocal identity is equivalent to
\[
 \boxed{
 H(w)=4^n w^nH\!\left(\frac1{16w}\right).
 }
\]
Thus the finite paired-root polynomial carries an exact functional relation inherited from the negative inversion.

\section{Claim boundary}

The following statements are proved:
\begin{itemize}
 \item $\lambda=4w$ converts $J(w)=1/(16w)$ to reciprocal inversion;
 \item the normalized paired set is reciprocal and conjugation invariant;
 \item its canonical monic polynomial has constant term $1$;
 \item $Q(x)=x^nQ(1/x)$;
 \item the coefficients of $Q$ are real and palindromic;
 \item the polynomial factors into reciprocal quadratic orbit factors and at most one exceptional $(x+1)$ factor;
 \item $Q(1)\neq0$ and $Q(-1)=0\iff F(-1/4)=0$.
\end{itemize}

The following are not proved:
\begin{itemize}
 \item nonemptiness of $P_J$;
 \item the zero or nonzero status of $F(-1/4)$ by a new analytic theorem;
 \item real-rootedness of $F$;
 \item PF-infinity;
 \item SOH-G003;
 \item the Riemann Hypothesis.
\end{itemize}
